\section{Saturated interval bundles and joint geometric cancellation}
\label{sec:saturated-bundles}

The existing primitive-annulus construction\cite{repo} applies when a bordered base
surface has a boundary circle that is not fully attached to the core.
We now remove that restriction under a different, explicitly checkable
condition.  An entire saturated bundle contributes one additional
two-handle relation.  We cancel these relations jointly with the original
core two-handles, using only the original core cocores.  The genera of the
bases can remain in binary throughout.

Here a \emph{saturated} bundle is an orientable total interval bundle
$P_i\to F_i$ over a connected compact surface with nonempty boundary, all
of whose vertical boundary annuli are attached to the core.  The base
$F_i$ need not be orientable.  We work with actual embedded annuli with
disjoint collars; an abstract cyclic word is insufficient to establish
this geometric hypothesis.  In the normal-cut application the source
parallelity construction and its ordered boundary observer supply these
annuli and collars.

The incoming hierarchy report\cite{hierarchy-incoming} already proves the
one-two-handle count used below. The new content of the relative lemma is
the exact old-cocore incidence formula and its implications for excluding
all unexpanded band columns from the raw search.

\subsection{One relative relation, with exact source incidences}

Let $U$ be a disjoint union of handlebodies given by zero- and one-handles.
Let $\alpha_1,\ldots,\alpha_f$ be the pairwise disjoint attaching curves of
the original core two-handles.  The core has no three-handles.  Fix a
spanning forest in the original zero/one-handle graph, and let
$\{D_h:h\in J\}$ be the cocore disk system corresponding to the remaining
one-handles.  Each saturated bundle is attached along annuli disjoint
from the $\alpha_r$.  Write $\beta_{ij}$ for an actual embedded core curve
of its $j$th attaching annulus, viewed on the $U$ side.

All curves are transverse to the cocore boundaries.  Define
\begin{equation}
 a_{rh}=|\alpha_r\cap\partial D_h|,\qquad
 a_{ih}^{\mathrm{sat}}
   =\sum_{j=1}^{b_i}|\beta_{ij}\cap\partial D_h|,
 \label{eq:saturated-incidences}
\end{equation}
where $b_i=|\pi_0(\partial F_i)|$.  These are unsigned geometric counts.
Oppositely oriented passages both contribute, and no free reduction is
performed.

\begin{lemma}[Relative bundle presentation with exact incidences]
\label{lem:saturated-relative-handles}
The union of $U$ with the saturated bundles has a relative handle
decomposition with, for each $P_i$, precisely
\[
 h_i=1-\chi(F_i)
\]
new one-handles and one new two-handle, and no new zero- or three-handles.
The new two-handle for $P_i$ has incidence
$a_{ih}^{\mathrm{sat}}$ with every old cocore $D_h$.  It traverses each
of its own new one-handles exactly twice and every other bundle's new
one-handles zero times.  Original core two-handles traverse all the new
one-handles zero times.
\end{lemma}

\begin{proof}
Choose pairwise disjoint properly embedded arcs in $F_i$ cutting it into
one disk.  Their number is $1-\chi(F_i)$: each cut increases Euler
characteristic by one.  Equivalently, $F_i$ is obtained from collars of
all its boundary circles by attaching $h_i$ bands and then one final
disk.  This relative surface handle decomposition exists for orientable
and nonorientable bordered surfaces alike.

Thicken this construction using the given interval bundle.  A boundary
collar can be absorbed into a collar of its annular attachment site in
$U$.  Each band thickens to a three-dimensional one-handle, and the
final disk, whose interval bundle is trivial, thickens to a two-handle.
This proves the asserted handle counts.

The boundary of the final disk follows each long side of each band once.
Its attaching curve therefore crosses every new band cocore exactly
twice.  A twisted band can change the signs of those crossings, but
cannot change this unsigned count.  The same boundary follows each
original boundary segment of $F_i$ once.  Those segments partition the
original boundary circles, so every original boundary circle is followed
once in total.  A whole boundary circle need not occur consecutively in
the attaching word.  Counting its segments avoids choices of basepoints
or conjugating paths.

For the exact counts on $U$, first straighten the annular attaching
collars about their chosen core curves. A relative Dehn twist of a
vertical annulus extends across a collar in the interval bundle: in
collar coordinates it has the form
$(\theta,r,t)\mapsto(\theta+2\pi k t\rho(r),r,t)$, where $\rho$ is one
at the vertical boundary and zero at the inner edge. Its effect on
the horizontal boundary is allowed because the conclusion here is
unmarked. Thus annular gluing parameters introduce no additional
restriction on the positions of the band feet.

Now narrow each annular attaching collar
about its chosen core curve $\beta_{ij}$, fixing that core curve.  In
transverse position, the intersections with the cocore boundaries then
consist of short spanning arcs, one for each counted intersection with
$\beta_{ij}$.  The finite band feet can be chosen in strips avoiding
all these arcs.  Consequently the old-boundary portions of the final
attaching curve contain every counted crossing exactly once, and the
new bands and their feet introduce no further crossing of an old
cocore.  This proves~\eqref{eq:saturated-incidences}.  The collar
changes are supported away from the original two-handle attaching
curves and the other marked interfaces.

Finally, every added band belongs to its own bundle.  No original
two-handle or other bundle's final disk uses it, giving the remaining
zero incidences.
\end{proof}

For a disk base, $h_i=0$, and the lemma is exactly the familiar absorption
of a fully attached disk bundle as a two-handle.  An annulus and a
M\"obius band both have $h_i=1$.  Their additional one-handle occurs twice
in the residual relation; its orientation data can differ, but the
incidence argument is the same.

\subsection{Why the additional columns can be omitted completely}

Let $A$ be the matrix obtained by stacking the $f$ original core rows and
one row $a_i^{\mathrm{sat}}$ per saturated bundle, keeping only the old
nonforest columns $J$.  In particular, a bundle contributes one row even
when $h_i$ is very large.

Extend the fixed original forest to a forest for the enlarged
zero/one-handle graph using only added band edges.  This is possible
because the original forest already spans each original core component.
Contraction deletes the selected added columns and leaves all old
nonforest incidences unchanged.  Let $N$ denote the full remaining
matrix, including the added band columns that were not contracted.

\begin{theorem}[Exact reduction to the aggregate source matrix]
\label{thm:saturated-even-columns}
The full relative matrix $N$ admits a complete sequence of positive
geometric singleton pivots if and only if $A$ does.  Thus the source
column-leaf algorithm decides complete cancellation in this selected
relative handle decomposition without constructing any added band
column.
\end{theorem}

\begin{proof}
Every entry in an added band column is initially even, by
Lemma~\ref{lem:saturated-relative-handles}.  Recall that a positive
pivot $(r,g)$, with $n_{rg}=1$, updates surviving entries by
\[
 n'_{ih}=n_{ih}+n_{ig}n_{rh}
 \qquad(i\ne r,\ h\ne g).
\]
If column $h$ has only even entries, its updated entries are even as
well.  Inductively, no added band column can ever contain a singleton.
Every possible pivot must therefore use an original core column.

Restricted to the original core columns, each update is precisely the
same update applied to $A$.  A complete sequence for $N$ consequently
gives one for $A$.  Conversely, a complete sequence for $A$ can be applied
to $N$ with the same pivot entries, deleting all its two-handle rows.
Theorem~\ref{thm:complete}, which equates positive singleton cancellation
with source column-leaf deletion, completes the argument.
\end{proof}

In fact, evenness is stronger than the algebra needs. A column whose
entries are all zero or at least two can never acquire an entry one under
a positive update: both summands are zero or at least two. Thus any such
column can be omitted when deciding full raw cancellation. The geometric
band construction gives the particularly transparent values zero and two.

The qualifier ``selected relative handle decomposition'' matters.
The theorem permits every positive pivot in that decomposition, not
arbitrary changes of disk system or arbitrary handle moves.  An
uncancelled stopping set is an obstruction to this rule, not an
obstruction to the underlying manifold being a handlebody.

\begin{corollary}[Joint core and saturated-bundle certificate]
\label{cor:saturated-handlebody}
If source column-leaf deletion removes all rows of $A$, then the core
with all saturated bundles attached is a disjoint union of
handlebodies.  The cancellation order may interleave original core
two-handles and saturated-bundle relations.
\end{corollary}

\begin{proof}
Every pivot has the geometric realization of a one/two-handle
cancellation.  For a source column leaf, no surviving attaching curve
meets the pivot cocore, so the positive update has no fill in any
surviving row.  Lemma~\ref{lem:saturated-relative-handles} authenticates
the geometric rows, and Theorem~\ref{thm:saturated-even-columns} shows
that the omitted columns hide no extra two-handle relation.  After all
rows have been removed, only zero- and one-handles remain.  Each
connected component is therefore a handlebody.
\end{proof}

\subsection{A primitive anchor explains the geometry}

The following special case makes the saturated-bundle step intuitive
without introducing its full relative matrix.
In this proposition, the remaining marked interfaces mean the other
attachment sites on $H$ and the displayed vertical annuli. No arbitrary
graph on the horizontal boundary of $P$ is included in that phrase.

\begin{proposition}[One primitive anchor releases a saturated bundle]
\label{prop:saturated-anchor}
Let $H$ be a handlebody and let $P\to F$ be an orientable interval bundle
over a connected bordered surface.  Choose one vertical annulus $A_0$.
Suppose its attachment annulus on $H$ has a properly embedded dual disk
$E$ meeting that annulus in one spanning arc and avoiding all other
marked interfaces.  After attaching $P$ to $H$ only along $A_0$, the
result is a handlebody of genus
\[
 g(H)+1-\chi(F)-1=g(H)-\chi(F).
\]
All remaining vertical annuli of $P$ form a simultaneous primitive
family in this handlebody, relative to the remaining marked interfaces.
\end{proposition}

\begin{proof}
For the first annular gluing, cut the annulus along a spanning arc to
obtain a disk interface. Restoring the annulus supplies a connecting
one-handle and a two-handle. The latter crosses the handle dual to $E$
exactly once, so that pair cancels geometrically. The pieces are then
joined by the connecting handle, giving the claimed genus. This is the
usual primitive-annulus argument\cite{schultens-weidmann}; the explicit
dual-disk construction is the hypothesis used here.

For every other boundary circle $C_j$ of $F$, choose a properly embedded
arc $\gamma_j$ joining it to $C_0$, with the arcs pairwise disjoint.
Such arcs can be chosen successively by joining distinct boundary
components; cutting along one leaves the base connected and merges
those two boundary components.  This construction does not require
orientability of $F$.

Each $p^{-1}(\gamma_j)$ is a rectangle with one edge on $A_0$ and one
edge on the annulus over $C_j$.  Choose pairwise disjoint parallel
copies $E_j$ of $E$.  A collar isotopy matches their spanning arcs in
$A_0$ with the corresponding rectangle edges: any finite disjoint
family of spanning arcs of an annulus is parallel.  Gluing each
rectangle to its $E_j$ along that edge gives an embedded proper disk
in $H\cup_{A_0}P$.  These disks are pairwise disjoint.  Each meets its
designated remaining annulus in one spanning arc and avoids all the
others, proving simultaneous primitivity.
\end{proof}

A singleton in an aggregate saturated row supplies exactly such an
anchor after the preceding source cancellations.  The relative-handle
proof is more general: anchors need not exist simultaneously in the
original core and may become available only after other rows disappear.

\subsection{Completing the census and its complexity}

After the joint cancellations, attach every unsaturated bundle using
the established spare-boundary primitive-annulus construction, and
attach the remaining proper disk interfaces.  These operations yield
handlebodies.  They can be delayed until after the joint cancellation
because their disjoint marked sites persist under the resulting
homeomorphisms.  No explicit coordinates for their transported images
are required for this unmarked topological conclusion.

Let $C$ be a connected component of the full attachment graph.  Its
Euler characteristic is
\begin{equation}
 \chi_C
 =\sum_{K_v\subset C}\chi(K_v)
  +\sum_{P_i\subset C}\chi(F_i)
  -d_C,
 \qquad g_C=1-\chi_C,
 \label{eq:saturated-genus}
\end{equation}
where $d_C$ is the number of proper disk interfaces in $C$; annuli have
Euler characteristic zero.  Equation~\eqref{eq:saturated-genus} computes
the genus only after the cancellation and gluing proof establishes that
$C$ is a handlebody.  It cannot be used as a stand-alone ball test.
An unported bordered bundle is an isolated handlebody of genus
$1-\chi(F_i)$ and retains its full binary multiplicity.

For the saturated subassembly alone, suppose a connected component
contains $k$ original zero/one-handlebody components, of genera $g_v$,
and $f_C$ original core two-handles.  The same count becomes
\[
 g_C=\sum_v g_v-f_C-k+1-\sum_i\chi(F_i).
\]
It is independent of how many added band edges were needed to join the
original components.

Let $r$ be the number of aggregate and original rows, $c$ the number of
old nonforest columns, and $I$ the number of nonzero source incidences.
After source reconstruction, the queue-based column-leaf algorithm uses
$O(I+r+c)$ structural operations; a heap for a deterministic least-label
order uses $O(I+(r+c)\log(r+c+1))$.  Counts can be saturated at two,
because only zero, one, and at least two are relevant.  The certificate
lists at most $r$ source row/column pairs and can be checked by one
incidence scan together with deletion times.  Base Euler
characteristics, genera, and isolated multiplicities are manipulated in
binary.  No operation is proportional to $\sum_i(1-\chi(F_i))$.

The complete normal-cut wrapper must additionally authenticate the
normal vector, compressed parallelity profiles, original handle charts,
and cyclic annular interface words.  Its total bound includes those
existing source routines and their binary arithmetic.  Closed base
surfaces and core three-handles remain outside this certificate.
Failure of the source cancellation test returns an inconclusive result.

This extension proves the unmarked homeomorphism type of more supplied
cuts.  It supplies neither an inherited boundary-pattern encoding for
the next hierarchy stage nor a bound on the number of candidate
surfaces or hierarchy resets.  Those are the remaining global issues
that a quasi-polynomial unknot-recognition argument must address.
